\section*{第 \ref{ch:b2:maxwell-equations} 章——麦克斯韦方程组}

\begin{solution}{exo:b2:maxwell-equations:1}
(a) $\operatorname{div} = 3$，$\operatorname{\vect{curl}} = \vect 0$：$r^2/2$ 的\omterm{def:b2:maxwell-equations:operators}{梯度}。
(b) $0$ 与 $(0, 0, 2)$：一个\omterm{def:b2:maxwell-equations:operators}{旋度}（$-\tfrac12(x^2 + y^2)\vect e_z$ 的\omterm{def:b2:maxwell-equations:operators}{旋度}），
不是\omterm{def:b2:maxwell-equations:operators}{梯度}。(c) $0$ 与 $\vect 0$：$xyz$ 的\omterm{def:b2:maxwell-equations:operators}{梯度}（同时也是\omterm{def:b2:maxwell-equations:operators}{旋度}）。(d)
$\frac1{r^2}\partial_r(r^2/r^2) = 0$，\omterm{def:b2:maxwell-equations:operators}{旋度} $\vect 0$：$-1/r$ 的\omterm{def:b2:maxwell-equations:operators}{梯度}。(e)
$0$ 与 $(0, -2x, 0)$：一个\omterm{def:b2:maxwell-equations:operators}{旋度}，不是\omterm{def:b2:maxwell-equations:operators}{梯度}。
\end{solution}

\begin{solution}{exo:b2:maxwell-equations:2}
(a) 高斯定理与法拉第定律：两边单位均为 \unit{V/m^2}（\omterm{def:b2:charges-currents-conduction:densities}{电荷密度}除以
$\varepsilon_0$，场的\omterm{def:b2:maxwell-equations:operators}{梯度}，场对时间）；安培定律：两边单位均为
\unit{T/m}。
(b) $1/\sqrt{8.85 \times 10^{-12} \times 1.257 \times 10^{-6}} = \qty{3.00e8}{m/s}$。(c) 比值
$\varepsilon_0\omega/\gamma$：铜为 $5 \times 10^{-20}$；\qty{1}{GHz} 下的海水为 $1 \times 10^{-2}$
（仍为导体，勉强如此）；\qty{50}{Hz} 下的玻璃为 $3 \times 10^3$：电介质，
\omterm{thm:b2:maxwell-equations:maxwell}{位移电流}占主导。
\end{solution}

\begin{solution}{exo:b2:maxwell-equations:3}
$\dd E/\dd t = I/\varepsilon_0\pi R^2 = \qty{1.8e12}{V.m^{-1}.s^{-1}}$; $j_d = I/\pi R^2 =
\qty{16}{A/m^2}$; $B(\qty{5}{cm}) = \mu_0Ir/2\pi R^2 = \qty{0.5}{\micro T}$; $B(\qty{20}{cm}) =
\mu_0I/2\pi r = \qty{0.5}{\micro T}$——恰为导线的磁场。
\end{solution}

\begin{solution}{exo:b2:maxwell-equations:4}
(a) $\rho = \varepsilon_0k$。(b) $\frac1{r^2}\partial_r(\rho_0r^3/3\varepsilon_0) = \rho_0/\varepsilon_0$；
球外由高斯定理，$E = \rho_0R^3/3\varepsilon_0r^2$。(c) $\operatorname{div}\vect B = 0$：成立；
$\operatorname{\vect{curl}}\vect B = -(B_0/a)\vect e_z$：$\vect j = -(B_0/\mu_0a)\vect e_z$，是
沿 $y$ 方向厚度均匀的\omterm{prop:b2:maxwell-equations:boundary}{面电流}。
\end{solution}

\begin{solution}{exo:b2:maxwell-equations:5}
(a) $\vect A = \tfrac12(B_yz - B_zy, B_zx - B_xz, B_xy - B_yx)$：$(\operatorname{\vect{curl}}
\vect A)_x = \tfrac12(B_x + B_x) = B_x$，等等；$\operatorname{div}\vect A = 0$。(b)
$\vect A' - \vect A = \tfrac12B_0(y, x, 0) = \operatorname{\vect{grad}}(\tfrac12B_0xy)$。(c)
$\frac1r\partial_r(Br^2/2) = B$；$\frac1r\partial_r(BR^2/2) = 0$；两者都等于 $BR/2$
在 $R$ 处。(d) $2\pi r\cdot BR^2/2r = \pi R^2B = \Phi$：螺线管外的一个回路，
其中 $\vect B = \vect 0$，仍感受到 $e = -\dd\Phi/\dd t$——通过 $\vect E =
-\partial_t\vect A$ 实现，而该处的矢势并不为零。
\end{solution}

\begin{solution}{exo:b2:maxwell-equations:6}
(a) $E = \sigma/\varepsilon_0 = \qty{1.1e5}{V/m}$，沿法向。(b) 两板之间为 $\sigma/\varepsilon_0$，板外
为零（每次跃变都增加 $\pm\sigma/\varepsilon_0$）。(c) $j_s = nI = \qty{1e4}{A/m}$，跃变
$\mu_0j_s = \qty{12.6}{mT} = B$ 在环内。(d) 环外的一个回路所包围的
净电流为零，故该处 $B = 0$；环内由安培环路定理得 $\mu_0NI/2\pi r$，
而绕组两侧的跃变 $\mu_0j_s = \mu_0NI/2\pi r$，二者一致。
\end{solution}

\begin{solution}{exo:b2:maxwell-equations:7}
(a) $V'' = -\rho_0/\varepsilon_0$：$V = -\rho_0x^2/2\varepsilon_0 + (U/d + \rho_0d/2\varepsilon_0)x$。(b)
$E = \rho_0x/\varepsilon_0 - U/d - \rho_0d/2\varepsilon_0$，在 $x_0 = d/2 + \varepsilon_0U/\rho_0d$ 处为零
——空间电荷较大时，零点靠近中点。(c) $\sigma(0) = \varepsilon_0E(0^+) =
-\varepsilon_0U/d - \rho_0d/2$，$\sigma(d) = -\varepsilon_0E(d^-) = \varepsilon_0U/d - \rho_0d/2$：两者合计
$-\rho_0d$，正好抵消空间电荷。(d) $\rho_0 = 0$：$V$ 为线性，
$E$ 均匀；$U = 0$：在 $d/2$ 处达到峰值的抛物线。
\end{solution}

\begin{solution}{exo:b2:maxwell-equations:8}
(a) \qty{6000}{km}，\qty{300}{m}，\qty{3}{m}，\qty{12.5}{cm}。(b) 当 $f < c/10L$ 时：
\qty{300}{MHz}；\qty{3}{GHz}；\qty{30}{Hz}。(c) $(2\pi \times 10^6 \times 0.05/3 \times 10^8)^2/4 =
3 \times 10^{-7}$。(d) 严格来说并不成立：在 \qty{50}{Hz} 下，长 \qty{1000}{km} 的电网
跨度仅为六分之一波长，各处相位不同，而长输电线路
应当按\omterm{prop:b2:dispersion-wave-packets:coax}{传输线}处理——不过每段局部电路仍是
似稳的。
\end{solution}

\begin{solution}{exo:b2:maxwell-equations:9}
$\operatorname{div}\operatorname{\vect{curl}}\vect B = 0$ 将要求 $\operatorname{div}\vect j =
0$，即 $\partial_t\rho = 0$ 处处成立。在极板（体积 $\pi R^2e$）中，
电流 $I$ 流入，无任何流出：$\langle\operatorname{div}\vect j\rangle =
-I/\pi R^2e = \qty{-1.6e5}{A/m^3}$, $\partial_t\rho = \qty{+1.6e5}{C.m^{-3}.s^{-1}}$.
\end{solution}

\begin{solution}{exo:b2:maxwell-equations:10}
(a) $2\pi rE_\theta = -\pi r^2\dot B$：内部 $E_\theta = -r\dot B/2$，
外部 $-R^2\dot B/2r$。(b) 其场线为闭合圆周：它不是\omterm{def:b2:maxwell-equations:operators}{梯度}，
它驱使电荷绕回路运动——是一种电动势场。(c) $e\cdot2\pi r_0
|E_\theta| = e\pi r_0^2\dot B$：$\pi \times 0.25 \times 10 = \qty{7.9}{eV}$ 每圈；$2.5 \times 10^6$
圈。(d) $p = eB(r_0)r_0$ 给出 $\dot p = er_0\dot B(r_0)$，而感生电场
给出 $\dot p = e|E_\theta| = e\dot\Phi/2\pi r_0$：$B(r_0) = \Phi/2\pi r_0^2 =
\tfrac12\langle B\rangle$.
\end{solution}

\begin{solution}{exo:b2:maxwell-equations:11}
(a) $\operatorname{\vect{curl}}\operatorname{\vect{curl}}\vect E = -\Delta\vect E$（因为 $\operatorname{div}
\vect E = 0$）$= -\partial_t\operatorname{\vect{curl}}\vect B = -\mu_0\varepsilon_0\partial_t^2\vect E$。(b)
对换角色后结论相同。(c) $c = 1/\sqrt{\mu_0\varepsilon_0} = \qty{3.0e8}{m/s}$。
(d) 韦伯与科尔劳施曾测得电荷的静电单位与电磁单位之比，
其值为 $3.1 \times 10^8$ m/s；菲佐曾测得
光速为 $3.15 \times 10^8$ m/s：“我们几乎
无法避免这样的推断：光是
同一介质中的横向波动。”
\end{solution}

\begin{solution}{exo:b2:maxwell-equations:12}
(a) $\operatorname{\vect{curl}}\operatorname{\vect{curl}}\vect B = -\Delta\vect B = \mu_0\gamma\operatorname{\vect{curl}}
\vect E = -\mu_0\gamma\partial_t\vect B$：$\Delta\vect B = \mu_0\gamma\partial_t\vect B$，即
扩散系数为 $1/\mu_0\gamma$ 的热传导方程形式。(b) $\tau \sim \mu_0\gamma d^2$：
极板为 \qty{7.5}{ms}；铁芯为 $4\pi \times 10^{-7} \times 5 \times 10^5 \times 9 \times 10^{12} = \qty{6e12}{s}$，约二十
万年。(c) $\mu_0\gamma\delta^2 \sim 2/\omega$：$\delta = \sqrt{2/\mu_0
\gamma\omega}$：在 \qty{50}{Hz} 时为 \qty{9}{mm}，在 \qty{1}{MHz} 时为 \qty{65}{\micro m}。(d) 铁芯的场
无法通过扩散在不到 $10^5$ 年内改变，
观测到的长期变化与倒转要求
导电流体的运动——即发电机。
\end{solution}

\begin{solution}{pb:b2:maxwell-equations:1}
\textbf{1.} $Q = (I_0/\omega)\sin\omega t$，$E = Q/\varepsilon_0\pi R^2$：振幅为 $I_0/\omega\varepsilon_0
\pi R^2 = \qty{5.7e5}{V/m}$.

\textbf{2.} $j_d = \varepsilon_0\partial_tE = I(t)/\pi R^2 = 32\cos\omega t$ A/m$^2$；其
穿过极板的通量为 $I(t)$。

\textbf{3.} $2\pi rB = \mu_0\varepsilon_0\pi r^2\partial_tE$：$B = \mu_0I(t)r/2\pi R^2$；在 $R$ 处：
$\mu_0I_0/2\pi R = \qty{2}{\micro T}$。

\textbf{4.} 对 $r > R$，全部\omterm{thm:b2:maxwell-equations:maxwell}{位移电流}都被包围在内：$B =
\mu_0I/2\pi r$，即导线的磁场；在 $R$ 处连续。

\textbf{5.} $(\operatorname{\vect{curl}}\vect E_1)_\theta = -\partial_rE_1 = -\partial_tB_\theta$，于是 $\partial_rE_1
= \mu_0\varepsilon_0(r/2)\ddot E$，且 $E_1(r) - E_1(0) = \mu_0\varepsilon_0r^2\ddot E/4$；代入
$\ddot E = -\omega^2E$：$|E_1(R) - E_1(0)|/|E| = (\omega R/c)^2/4 = 1 \times 10^{-6}$。

\textbf{6.} 很好。$10\%$ 时需 $(\omega R/c)^2 = 0.4$，$f \approx \qty{300}{MHz}$——
此时为谐振腔，不再是电容器。

\textbf{7.} $W_e = \tfrac12\varepsilon_0E_0^2\pi R^2d = \qty{4.5e-5}{J}$；$W_m = \int_0^R(\mu_0I_0r/
2\pi R^2)^2/2\mu_0\cdot2\pi rd\,\dd r = \mu_0I_0^2d/16\pi = \qty{2.5e-11}{J}$；比值 $6 \times
10^{-7} \approx (\omega R/c)^2/8$：二者为同一小参数。

\textbf{8.} $\oint\vect E\cdot\dd\vect l = 2\pi r_0E_\theta = -\dd\Phi/\dd t$。

\textbf{9.} $\dd p/\dd t = e|E_\theta| = (e/2\pi r_0)|\dd\Phi/\dd t|$；每转一圈
有 $2\pi r_0\cdot e|E_\theta| = e|\dd\Phi/\dd t|$。

\textbf{10.} 由 $p = eB(r_0)r_0$ 得 $\dot p = er_0\dot B(r_0)$；与第 9 题联立：
$B(r_0) = \Phi/2\pi r_0^2 = \tfrac12\langle B\rangle$。

\textbf{11.} 均匀磁场使 $B(r_0) = \langle B\rangle$，为所需值的两倍：
磁力增长快于动量，轨道
收缩。磁极的形状使磁场在靠近
轴线处强，在轨道上则只有一半。

\textbf{12.} $p = eBr_0 = \qty{3.2e-20}{kg.m/s}$；$E = pc = \qty{9.6e-12}{J} = \qty{60}{MeV}$；
$\Phi = 2\pi r_0^2B = \qty{0.63}{Wb}$ 在 \qty{5}{ms} 内：每圈 \qty{126}{eV}，$4.8 \times
10^5$ 圈，\qty{1500}{km}——以 $c$ 计，\qty{5}{ms} 内所走距离相同：一致。

\textbf{13.} 能量来自磁铁的电源，经
法拉第项交给电子：变化的磁通量产生加速电场。

\textbf{14.} 数十 MeV 的电子在重靶中制动，产生
穿透力强的 X 射线；紧凑型直线加速器，剂量
率更高，取代了它。

\textbf{15.} $\operatorname{div}\operatorname{\vect{curl}}\vect A = 0$；$\operatorname{\vect{curl}}(-\operatorname{\vect{grad}}
V - \partial_t\vect A) = -\partial_t\operatorname{\vect{curl}}\vect A = -\partial_t\vect B$。

\textbf{16.} $\frac1{r^2}(r^2V')' = -\rho_0/\varepsilon_0$：$V = \rho_0(3a^2 - r^2)/6\varepsilon_0$
内部，外部为 $\rho_0a^3/3\varepsilon_0r = Q/4\pi\varepsilon_0r$；$E = \rho_0r/3\varepsilon_0$ 与
$Q/4\pi\varepsilon_0r^2$。

\textbf{17.} $\tfrac12\int\rho V\,\dd\tau = 2\pi\rho_0^2(a^5/6 - a^5/30)/\varepsilon_0 = 4\pi\rho_0^2a^5/
15\varepsilon_0 = 3Q^2/20\pi\varepsilon_0a$；铀：\qty{1.6e-10}{J} $\approx \qty{1000}{MeV}$（裂变
释放的 \qty{200}{MeV}，即两个较小的球所省的能量）。

\textbf{18.} $V = -E(t)z$（另加一常数）。设 $\vect A = A_z(r, t)\vect e_z$，
由 $B_\theta = -\partial_rA_z = \mu_0Ir/2\pi R^2$ 得 $A_z = -\mu_0I(t)r^2/4\pi R^2$。

\textbf{19.} $-\partial_tA_z = \mu_0\dot Ir^2/4\pi R^2 = \mu_0\varepsilon_0r^2\ddot E/4$（因 $\dot I =
\varepsilon_0\pi R^2\ddot E$）：即第 5 题中的修正，差一个常数。

\textbf{20.} $\operatorname{div}\vect A = \partial_zA_z = 0$，因 $A_z$
仅依赖于 $r$；任意 $\vect A + \operatorname{\vect{grad}}\chi$ 亦然。

\textbf{21.} $\sigma = \varepsilon_0E_0 = \qty{5}{\micro C/m^2}$；金属内部 $\vect E =
\vect 0$，故跃变 $\sigma/\varepsilon_0$ 恰为间隙场。

\textbf{22.} 电流 $I(t)$ 从轴心向极板供给；超过 $r$ 的、
仍待输送的部分为 $I(1 - r^2/R^2)$，故 $j_s = I(1 -
r^2/R^2)/2\pi r$。在外表面，由安培环路定理得导线的磁场 $\mu_0I/
2\pi r$；在内表面为 $\mu_0Ir/2\pi R^2$；二者之差 $\mu_0I(1 -
r^2/R^2)/2\pi r = \mu_0j_s$，即边界跃变。

\textbf{23.} 极板：$R/\lambda = 3 \times 10^{-4}$，电似稳。回路：$7 \times 10^{-4}$，
磁似稳。在 \qty{1}{GHz} 时：$R/\lambda = 0.3$——既是谐振腔，也是天线。

\textbf{24.} 对于米级的环，$\lambda = \qty{3000}{km}$：\omterm{thm:b2:maxwell-equations:maxwell}{位移
电流}完全可以忽略；起主要作用的电场
是法拉第感应电场，磁似稳完整地保留了它。

\textbf{25.} I：安培–麦克斯韦（\omterm{thm:b2:maxwell-equations:maxwell}{位移电流}）、法拉第（修正项）、
能量密度。II：法拉第积分、洛伦兹力。
III：势、泊松（高斯）。IV：$\vect E$
和 $\vect B$ 的边界跃变；由尺寸与波长之比判定所属区域。
\end{solution}
